% !TEX program = xelatex
\documentclass[xcolor={dvipsnames,table,hyperref}]{beamer}

% pick one. "hide notes" is default
% \setbeameroption{hide notes}
\setbeameroption{show notes}
% \setbeameroption{show notes on second screen = left}
% \setbeameroption{show only notes}
% \setbeameroption{show only slides with notes}

\usepackage{fontspec}
\usepackage{fontsetup}

% \usepackage{FiraSans}
% Use CM Bright for math calligraphic instead of Fira
\DeclareMathAlphabet{\mathcal}{OMS}{cmbrs}{m}{n}
% \setmonofont{Fira Mono}

\usepackage{polyglossia}
\setdefaultlanguage{english}

\usepackage{xltxtra}
\usepackage[backend=biber,style=alphabetic,
doi=true, url=true, isbn=false]{biblatex}

% When Jeffrey Goldberg builds this, he can use his bibliography database
% But he should include in the repository a .bib file that contains what
% is needed for these slides.
% He can do that by first generating the .bcf using his resources, but then
% using
%  biber --output_format=bibtex --output_resolve modern-auth.bcf
%  rm modern-auth.bcf
% to create a modern-auth_biber.bib file to be included.
%
\IfFileExists{./modern-auth_biber.bib}{%
  \addbibresource{modern-auth_biber.bib}
}{%
  \addbibresource{crypto.bib}
  \addbibresource{statistics.bib}
}

\renewcommand*{\bibfont}{\tiny}

% Beamer (metropolis) style setup
%\usepackage{FiraSans}
%\usefonttheme{professionalfonts}
\usetheme{moloch}
\definecolor{BitsBlue2}{HTML/hsb}{0A56BF/0.597,0.95,0.75} % guess % H=215
\molochset{titleformat = smallcaps,
  titleformat plain = regular,
  subsectionpage = progressbar,
block = fill}
% \setbeamercolor{palette primary}{bg=BitsBlue2}
\molochcolors{
  light/frametitle bg=BitsBlue2
}
\usepackage{appendixnumberbeamer}

\setbeamertemplate{theorems}[numbered]

\usepackage{Common/stNotes}
%\setbeamertemplate{frame footer}{\insertshortauthor}

\usepackage{minted}
\usemintedstyle{xcode}

\colorlet{Ucolor}{green!25}
\colorlet{Zcolor}{yellow!50}
\newcolumntype{U}{>{\columncolor{Ucolor}[1.0\tabcolsep]}r}
\newcolumntype{Z}{>{\columncolor{Zcolor}[1.0\tabcolsep]}r}
\newcommand{\rcu}{\rowcolor{Ucolor}}
\newcommand{\rcz}{\rowcolor{Zcolor}}

\author[J.~Goldberg]
{Jeffrey Goldberg\texorpdfstring{\\ \texttt{jeffrey@goldmark.org}}{}}
\title{Building Modern Authentication in the Before Time (2015)}
\subtitle{Putting theory into practices is hard}
\date{September 1, 2026\\
\footnotesize Last revised: \today }
% \institute{Not yet institutionalized}


\hypersetup{colorlinks=true, allcolors=black, urlcolor=magenta}
\providecommand*{\reporoot}{https://github.com/jpgoldberg/modern-auth-talk}
\providecommand*{\sourceroot}{\reporoot/blob/main/src}
\begin{document}
\maketitle

\begin{frame}{Who am I?}
  \begin{itemize}
    \item Jeffrey Goldberg \url{https://jeffrey.goldmark.org}
    \item Of the generation of (dropout) academics who ended up as sysadmins.
      \note[item]{Has never taken a programming course, but have taught some.}
    \item Tried (and failed) to solve the password problem in the 1990s.
    \item Got into information security and cryptography
    \item 2010–2023: Worked for 1Password, a password manager.
    \item Not authorized to speak on behalf of 1Password.
      \note[item]{I am much more than that, but trying to keep things relevant and to a single slide.}
  \end{itemize}

\end{frame}

\begin{frame}{1Password in 2015}
  \begin{itemize}
    \item A successful password manager from home users.
    \item Very strong reputation for privacy as well as security and usability.
      \note[item]{E.g., among the first (2012) to encrypt URLs and item titles and to use authenticated encryption (Encrypt-then-MAC).}
    \item Did not host customer data in any form.
    \item Hooks for users to sync their data using Dropbox and other third party solutions.
  \end{itemize}
\end{frame}

\begin{frame}{Limations}
  \begin{itemize}
    \item Maintaining sync mechanisms with different third party services getting harder.
    \item Sync systems that worked didn't have all the security requirements we wanted.
    \item Users could share subsets of their data with other members of families or teams.
  \end{itemize}
\end{frame}

%%%%%%%%%%%%%%%%%%%%%%%%%%%%%%%%%%%
%%%%%%%%%%%%%%%%%%%%%%%%%%%%%%%%%%%
\end{document}


\section{Resources}

\begin{frame}{Resources}
  \begin{itemize}
    \item \href{\reporoot/public/s/rsa.pdf}{These slides}
    \item Exceedingly \href{\reporoot/public/s/citm-rsa.pdf}{extensive notes} on the Cat in the Middle computation. Those notes include explanations of some of the algorithms discussed.
    \item \href{\sourceroot/}{Sources}
  \end{itemize}

\end{frame}

\begin{frame}[t,allowframebreaks]
  \frametitle{References}
  \printbibliography[heading=none]
\end{frame}

%%%%%%%%%%%%%%%%%%%%%%%%%%%%%%%%%%%
%%%%%%%%%%%%%%%%%%%%%%%%%%%%%%%%%%%
\end{document}
